\documentclass[11pt,a4paper]{article}
\usepackage[margin=2.5cm]{geometry}
\usepackage[T1]{fontenc}
\usepackage{booktabs}
\usepackage{xcolor}
\usepackage{amsmath}
\usepackage{listings}
\usepackage[hidelinks]{hyperref}

\emergencystretch=2em
\definecolor{codebg}{gray}{0.95}
\lstset{language=C++, basicstyle=\ttfamily\small, keywordstyle=\color{blue!70!black},
  commentstyle=\color{green!50!black}\itshape, backgroundcolor=\color{codebg},
  frame=single, rulecolor=\color{black!20}, breaklines=true, showstringspaces=false,
  tabsize=2, columns=flexible}

\title{ATmega32A Driver Course \\ \Large Module 6 --- UART (Serial Communication)}
\author{Ali Sahafi}
\date{}

\begin{document}
\maketitle

\section{Theory}

The UART (\emph{Universal Asynchronous Receiver/Transmitter}) sends bytes
over a single wire per direction: \textbf{TXD} (PD1) transmits,
\textbf{RXD} (PD0) receives. ``Asynchronous'' means there is no clock wire
--- both sides must agree on the \textbf{baud rate} (bits per second) in
advance. Each byte travels as a \emph{frame}: 1 start bit, 8 data bits, 1
stop bit (``8N1''), so at 9600 baud you get at most 960 bytes/s.

The hardware generates the baud rate by dividing the CPU clock. This driver
always enables \textbf{double speed mode} (\texttt{U2X}), which halves the
divider and reduces rounding error:

\begin{equation*}
\text{UBRR} = \frac{F_{\text{CPU}}}{8 \cdot \text{baud}} - 1
\end{equation*}

This is why a correct \texttt{F\_CPU} is critical: if it doesn't match the
real clock, the baud rate is wrong and the terminal shows garbage.

\section{Registers}

\begin{center}
\small
\begin{tabular}{lll}
\toprule
\textbf{Register} & \textbf{Role} & \textbf{Bits used by the driver} \\
\midrule
\texttt{UBRRH/UBRRL} & Baud rate divider & 12-bit UBRR value \\
\texttt{UCSRA} & Status & \texttt{U2X}, \texttt{UDRE} (TX ready), \texttt{RXC} (byte received) \\
\texttt{UCSRB} & Enable & \texttt{RXEN}, \texttt{TXEN}, \texttt{RXCIE} (RX interrupt) \\
\texttt{UCSRC} & Frame format & \texttt{URSEL}, \texttt{UCSZ1:0} = 8 data bits \\
\texttt{UDR}   & Data & write to send, read to receive \\
\bottomrule
\end{tabular}
\end{center}

\section{API reference}

\begin{center}
\small
\begin{tabular}{ll}
\toprule
\textbf{Method} & \textbf{Description} \\
\midrule
\texttt{UART.init(baud, interrupt)} & Configure baud rate, optional RX interrupt \\
\texttt{UART.transmit(byte)} & Send raw byte \\
\texttt{UART.print(value)} & Send string, char, int, uint or float (overloaded) \\
\texttt{UART.print(float, decimals)} & Float with chosen decimal places (default 2) \\
\texttt{UART.receive()} & Blocking receive --- waits for incoming byte \\
\texttt{UART.attachRxInterrupt(callback)} & Set RX callback: \texttt{void cb(uint8\_t byte)} \\
\bottomrule
\end{tabular}
\end{center}

Receiving can be \textbf{blocking} (\texttt{receive()} halts until a byte
arrives) or \textbf{interrupt-driven} (pass \texttt{true} to \texttt{init}
and attach a callback --- bytes then arrive ``by themselves'' while your main
loop does other work; remember \texttt{sei()}).

\section{Example: serial echo + uptime printer}

\paragraph{Wiring} (a USB--serial adapter, e.g. CP2102 or FT232):
\begin{itemize}
  \item \textbf{PD0 (RXD)} $\leftarrow$ adapter \textbf{TX}
  \item \textbf{PD1 (TXD)} $\rightarrow$ adapter \textbf{RX}
  \item GND $\leftrightarrow$ GND (common ground is mandatory!)
\end{itemize}
Open a serial terminal at \textbf{9600 baud, 8N1}.

\paragraph{Build \& flash.} \texttt{make uart}

\lstinputlisting{example.cpp}

\section{Exercises}

\begin{enumerate}
  \item Turn the echo into a \emph{command interpreter}: \texttt{'1'} switches
        an LED on, \texttt{'0'} switches it off, anything else prints
        \texttt{"?"}.
  \item Compute UBRR for 9600 baud at $F_{\text{CPU}} = 1$\,MHz (the factory
        default clock!) with U2X enabled. How large is the rounding error, and
        what do you expect to see in the terminal?
  \item Combine with Module 5: print the ADC value of the potentiometer once
        per second, and use a serial-plotter tool to graph it live.
\end{enumerate}

\end{document}
